%======================================================================
% CAPÍTULO 4: ARQUITECTURA DEL ECOSISTEMA
%======================================================================
\chapter{Arquitectura del Ecosistema}

\section{Visi\'on General}

El ecosistema adopta una arquitectura de microservicios donde cada componente cumple funciones espec\'ificas operando de manera coordinada.

\begin{table}[H]
\centering
\caption{Componentes del ecosistema y tecnolog\'{\i}as}
\begin{tabular}{|l|l|l|l|}
\hline
\textbf{Componente} & \textbf{Tecnolog\'{\i}a} & \textbf{Funci\'on} & \textbf{Puerto} \\ \hline
Sandra Server & Go 1.24.3 & Middleware/ESB & 443, 8443 \\ \hline
Sandra UI Consola & Angular 16+ & Panel Admin & 4200 \\ \hline
Sandra Desktop Container & Tauri 2.0 & Escritorio & - \\ \hline
Sandra Sentinel & Rust 2021 & Motor C\'alculo & 50051 \\ \hline
\end{tabular}
\end{table}

\section{Componentes Principales}

\subsection{Sandra Server}
N\'ucleo del ecosistema que proporciona servicios de integraci\'on, comunicaci\'on y seguridad. Implementa:
\begin{itemize}
\item WebSocket Hub con sharding (256 fragmentos)
\item API REST completa
\item M\'odulo de seguridad (JWT, MFA TOTP)
\item Sistema de mensajer\'{\i}a en tiempo real
\item Utilidades (PDF, QR, cifrado, compresi\'on, cache)
\end{itemize}

\subsection{Sandra UI Consola}
Panel de administraci\'on desarrollado en Angular 16+ con componentes standalone.

\subsection{Sandra Desktop Container}
Combinaci\'on de Rust con Angular mediante Tauri 2.0, funcionando como orquestador de entornos seguros.

\subsection{Sandra Sentinel}
N\'ucleo de procesamiento de alto rendimiento en Rust, especializado en c\'alculo de n\'ominas masivas.

\section{Modelo de Comunicaci\'on}

Los componentes se comunican mediante protocolos especializados:

\begin{itemize}
\item \textbf{REST/HTTPS}: Operaciones transaccionales
\item \textbf{WebSockets}: Tiempo real
\item \textbf{gRPC}: Transferencia eficiente de datos
\end{itemize}

Todas las comunicaciones est\'an cifradas mediante TLS 1.3, y los tokens JWT aseguran autenticaci\'on y autorizaci\'on en cada solicitud.

\subsection{Puertos y Protocolos}

\begin{table}[H]
\centering
\caption{Puertos y protocolos de comunicaci\'on}
\begin{tabular}{|l|l|l|}
\hline
\textbf{Servicio} & \textbf{Puerto} & \textbf{Protocolo} \\ \hline
API HTTPS & 443 & REST/HTTPS \\ \hline
API Desarrollo & 8443 & REST/HTTPS \\ \hline
WebSocket & 8080 & WSS \\ \hline
gRPC & 50051 & HTTP/2 \\ \hline
Consola Web & 4200 & HTTP \\ \hline
\end{tabular}
\end{table}

\newpage
